\section{Metodolog\'ia}

El presente estudio emplea una metodolog\'ia de c\'alculo anal\'itico complementada con simulaci\'on num\'erica. Los c\'alculos se desarrollan en el lenguaje de programaci\'on Python utilizando las bibliotecas cient\'ificas NumPy y SciPy, y se presentan en este informe con el procedimiento paso a paso para que el lector pueda verificar cada resultado de forma independiente. A continuaci\'on se describe la metodolog\'ia para el an\'alisis de transferencia de calor del sistema de media ca\~na.

\subsection{Determinaci\'on del Coeficiente Global de Transferencia de Calor}
\label{sec:metodologia_U}

El coeficiente global de transferencia de calor $U$ del sistema de media ca\~na se determina mediante el modelo de resistencias t\'ermicas en serie. Para una pared cil\'indrica plana, aproximaci\'on v\'alida dado que el espesor de l\'amina del tanque es peque\~no comparado con su di\'ametro, la expresi\'on general es:

\begin{equation}
\frac{1}{U} = \frac{1}{h_i} + \frac{e_{w}}{k_{w}} + \frac{1}{h_o}
\label{eq:U_global}
\end{equation}

donde $h_i$ es el coeficiente de convecci\'on del lado del agua caliente en el interior de la media ca\~na, $e_{w}$ es el espesor de la pared del tanque, $k_{w}$ es la conductividad t\'ermica del acero inoxidable SS316L, y $h_o$ es el coeficiente de convecci\'on del lado de la glucosa por convecci\'on natural.

El coeficiente del lado del agua se calcula mediante la correlaci\'on de Sieder-Tate para flujo turbulento en conductos no circulares, empleando el di\'ametro hidr\'aulico del canal rectangular:

\begin{equation}
D_h = \frac{4 \cdot A_c}{P_m} = \frac{4 \cdot a \cdot b}{2(a+b)}
\label{eq:Dh}
\end{equation}

donde $a$ y $b$ son las dimensiones internas del perfil rectangular, $A_c$ es el \'area de la secci\'on transversal de flujo y $P_m$ es el per\'imetro mojado. La correlaci\'on de Sieder-Tate se expresa como:

\begin{equation}
\text{Nu} = 0,027 \cdot \text{Re}^{0,8} \cdot \text{Pr}^{1/3} \cdot \left(\frac{\mu}{\mu_w}\right)^{0,14}
\label{eq:sieder_tate}
\end{equation}

v\'alida para $\text{Re} > 10{.}000$, $0,7 < \text{Pr} < 16{.}700$ y $L/D_h > 10$.

El coeficiente del lado de la glucosa se obtiene mediante la correlaci\'on de Churchill y Chu para convecci\'on natural sobre una superficie plana:

\begin{equation}
\text{Nu}_L = \left[0,825 + \frac{0,387 \cdot \text{Ra}_L^{1/6}}{\left(1 + (0,492/\text{Pr})^{9/16}\right)^{8/27}}\right]^2
\label{eq:churchill_chu}
\end{equation}

donde $\text{Ra}_L = \text{Gr}_L \cdot \text{Pr}$ es el n\'umero de Rayleigh, evaluado con la longitud caracter\'istica de la superficie de contacto y las propiedades de la glucosa a la temperatura de pel\'icula. La correlaci\'on fue desarrollada originalmente para placas verticales isotermas; su aplicaci\'on al fondo toriesf\'erico representa una aproximaci\'on conservadora, dado que la geometr\'ia curva calentada desde abajo favorece la convecci\'on natural respecto a una placa plana vertical.

\subsection{Modelo de Calentamiento Transitorio}

El calentamiento de la glucosa en el tanque se modela mediante un balance de energ\'ia transitorio. Para el caso de glucosa en reposo, sin descarga, la ecuaci\'on diferencial gobernante es:

\begin{equation}
m_{g} \cdot C_{p,g} \cdot \frac{dT_g}{dt} = U \cdot A \cdot \Delta T_{lm}
\label{eq:balance_transitorio}
\end{equation}

donde $m_g$ es la masa de glucosa en el tanque, $C_{p,g}$ es el calor espec\'ifico de la glucosa, $T_g$ es la temperatura media de la glucosa, $A$ es el \'area de contacto de la media ca\~na (14~m\textsuperscript{2}) y $\Delta T_{lm}$ es la diferencia de temperatura media logar\'itmica entre el agua de calentamiento y la glucosa:

\begin{equation}
\Delta T_{lm} = \frac{(T_{w,in} - T_g) - (T_{w,out} - T_g)}{\ln\left(\frac{T_{w,in} - T_g}{T_{w,out} - T_g}\right)}
\label{eq:LMTD}
\end{equation}

La ecuaci\'on se integra num\'ericamente mediante el m\'etodo de Runge-Kutta de cuarto orden con paso de tiempo adaptativo, actualizando las propiedades dependientes de la temperatura en cada paso.

Para los escenarios con descarga simult\'anea a carrotanque, el balance de energ\'ia se modifica para incluir el t\'ermino de salida de masa:

\begin{equation}
m_g \cdot C_{p,g} \cdot \frac{dT_g}{dt} = U \cdot A \cdot \Delta T_{lm} - \dot{m}_{out} \cdot C_{p,g} \cdot (T_g - T_{ref})
\label{eq:balance_descarga}
\end{equation}

donde $\dot{m}_{out}$ es el flujo m\'asico de descarga al carrotanque y $T_{ref}$ es la temperatura de referencia. Simult\'aneamente, la masa de glucosa var\'ia seg\'un:

\begin{equation}
\frac{dm_g}{dt} = -\dot{m}_{out}
\label{eq:balance_masa}
\end{equation}

El sistema de ecuaciones diferenciales ordinarias acopladas se resuelve con el m\'etodo RK45 adaptativo.

\subsection{Determinaci\'on del \'Area de Transferencia Requerida}
\label{sec:area_continuo}

Para evaluar la capacidad del sistema de chaqueta en modo de calentamiento continuo, se desarroll\'o un modelo de dimensionamiento basado en el m\'etodo de la diferencia de temperatura media logar\'itmica con integraci\'on num\'erica del coeficiente global variable. El calor requerido por unidad de tiempo para calentar un flujo m\'asico de glucosa $\dot{m}$ desde una temperatura de entrada $T_{g,in}$ hasta una temperatura de salida $T_{g,out}$ se expresa como:

\begin{equation}
\dot{Q} = \dot{m} \cdot C_{p,prom} \cdot (T_{g,out} - T_{g,in})
\label{eq:calor_continuo}
\end{equation}

donde $C_{p,prom}$ se eval\'ua a la temperatura media del rango de calentamiento. La diferencia de temperatura media logar\'itmica, considerando el agua a temperatura pr\'acticamente constante $T_w$ a lo largo de la chaqueta, se determina mediante:

\begin{equation}
\Delta T_{lm} = \frac{(T_w - T_{g,in}) - (T_w - T_{g,out})}{\ln\!\left(\dfrac{T_w - T_{g,in}}{T_w - T_{g,out}}\right)}
\label{eq:LMTD_continuo}
\end{equation}

Dado que $U$ var\'ia de manera significativa con la temperatura de la glucosa, el coeficiente global efectivo se obtiene mediante integraci\'on trapezoidal de la resistencia espec\'ifica sobre el rango de temperatura:

\begin{equation}
\frac{1}{U_{ef}} = \frac{1}{\Delta T_{g}} \int_{T_{g,in}}^{T_{g,out}} \frac{1}{U(T)} \, dT
\label{eq:U_efectivo}
\end{equation}

El \'area de transferencia requerida para alcanzar la temperatura de salida especificada resulta entonces:

\begin{equation}
A_{req} = \frac{\dot{Q}}{U_{ef} \cdot \Delta T_{lm}}
\label{eq:area_requerida}
\end{equation}

Se analizan dos rangos de temperatura de la glucosa: calentamiento completo desde 25~\textdegree C hasta 57~\textdegree C y calentamiento de mantenimiento desde 54~\textdegree C hasta 57~\textdegree C o desde 55~\textdegree C hasta 57~\textdegree C.

\subsection{Ciclo de Descargas a Carrotanque}

El ciclo operativo oficial consiste en cinco descargas diarias de 24~toneladas cada una. Cada descarga tiene una duraci\'on de 2~horas, lo que implica un flujo de descarga de 12~ton/h. El intervalo entre inicios de descarga es de 4,8~horas, de modo que el tiempo disponible para calentamiento entre el fin de una descarga y el inicio de la siguiente es de 2,8~horas. Durante las descargas, el calentamiento contin\'ua operando y la glucosa de alimentaci\'on ingresa al tanque a la temperatura de suministro. Entre descargas, no hay flujo de entrada ni salida y el sistema recupera la temperatura del inventario.

El flujo medio del ciclo resulta:

\begin{equation}
\dot{m}_{promedio} = \frac{5 \times 24{.}000~\text{kg}}{24~\text{h}} = 5{.}000~\text{kg/h}
\end{equation}

El ciclo se simula resolviendo las ecuaciones de balance transitorio con las condiciones de flujo variable: durante la descarga se impone $\dot{m}_{out} = 12{.}000$~kg/h y entre descargas $\dot{m}_{out} = 0$.

\subsection{Espesor de Aislamiento T\'ermico}

El espesor \'optimo de aislamiento de lana mineral se determina mediante un modelo de resistencias t\'ermicas en serie desde el agua caliente en la media ca\~na hasta el aire ambiente. La p\'erdida de calor por unidad de \'area se calcula como:

\begin{equation}
q = \frac{T_w - T_{amb}}{R_{conv,i} + R_{perfil} + R_{aisl} + R_{conv,e}}
\label{eq:aislamiento}
\end{equation}

donde $R_{conv,i} = 1/h_i$ es la resistencia de convecci\'on interna, $R_{perfil} = t_p/k_{SS}$ es la resistencia de conducci\'on de la l\'amina del perfil, $R_{aisl} = \delta/k_{aisl}$ es la resistencia del aislamiento de lana mineral y $R_{conv,e} = 1/h_e$ es la resistencia de convecci\'on natural en el aire exterior. Se eval\'ua el espesor m\'inimo requerido para que la temperatura superficial exterior no exceda 60~\textdegree C y se comparan las p\'erdidas de calor para espesores comerciales est\'andar.

\subsection{Validaci\'on Mediante Simulaci\'on CFD}

La simulaci\'on de din\'amica de fluidos computacional (CFD) se realiz\'o en COMSOL Multiphysics~6.4 para complementar los c\'alculos anal\'iticos. El objetivo fue obtener un valor representativo del coeficiente global de transferencia de calor $U$ bajo las condiciones del Escenario~3 y compararlo con el modelo anal\'itico implementado en Python. El dominio num\'erico incluye el volumen de agua caliente dentro de la chaqueta de media ca\~na en espiral, la l\'amina de acero inoxidable SS316L del fondo toriesf\'erico, y una regi\'on del volumen de glucosa Globe~1130 adyacente a la pared calentada. La geometr\'ia reproduce la configuraci\'on final de doble entrada y doble salida de agua caliente documentada en los planos mec\'anicos.

El modelo resuelve de forma acoplada las ecuaciones de conservaci\'on de masa, momento y energ\'ia. Para el flujo de agua en la media ca\~na, la ecuaci\'on de continuidad en r\'egimen estacionario e incompresible es:

\begin{equation}
\nabla \cdot \vec{v} = 0
\label{eq:continuidad_cfd}
\end{equation}

donde $\vec{v}$ es el vector velocidad del fluido. La ecuaci\'on de Navier-Stokes para la conservaci\'on del momento se expresa como:

\begin{equation}
\rho \left( \vec{v} \cdot \nabla \right) \vec{v} = -\nabla p + \nabla \cdot \left[ \mu \left( \nabla \vec{v} + (\nabla \vec{v})^T \right) \right]
\label{eq:navier_stokes_cfd}
\end{equation}

donde $\rho$ es la densidad del agua, $p$ es la presi\'on y $\mu$ es la viscosidad din\'amica. El t\'ermino convectivo se considera estacionario dado que el tiempo de residencia del agua en la chaqueta es varios \'ordes de magnitud menor que la constante de tiempo t\'ermica del tanque; la validaci\'on del coeficiente $U$ se realiza por tanto en condici\'on pseudo-estacionaria de flujo.

La ecuaci\'on de conservaci\'on de la energ\'a en el dominio fluido, despreciando el trabajo de disipaci\'on viscosa y la compresibilidad, es:

\begin{equation}
\rho C_p \left( \vec{v} \cdot \nabla T \right) = \nabla \cdot (k \nabla T)
\label{eq:energia_fluido_cfd}
\end{equation}

donde $C_p$ es el calor espec\'ifico, $T$ es la temperatura y $k$ es la conductividad t\'ermica del agua. Para la regi\'on s\'olida del fondo toriesf\'erico, la ecuaci\'on de conducci\'on de calor es:

\begin{equation}
\nabla \cdot (k_s \nabla T_s) = 0
\label{eq:conduccion_solido_cfd}
\end{equation}

donde $k_s$ y $T_s$ corresponden a la conductividad t\'ermica y la temperatura del acero SS316L, respectivamente.

En el dominio de glucosa se resuelve la ecuaci\'on de energ\'ia con el t\'ermino convectivo debido a la convecci\'on natural inducida por diferencias de densidad:

\begin{equation}
\rho_g C_{p,g} \left( \vec{v}_g \cdot \nabla T_g \right) = \nabla \cdot (k_g \nabla T_g)
\label{eq:energia_glucosa_cfd}
\end{equation}

donde el campo de velocidades $\vec{v}_g$ satisface la ecuaci\'on de momentum acoplada con el t\'ermino de flotaci\'on de Boussinesq, $\rho_g \vec{g} \beta (T_g - T_{ref})$, siendo $\beta$ el coeficiente de expansi\'on t\'ermica y $T_{ref}$ la temperatura de referencia del fluido en reposo. El movimiento natural permanece en r\'egimen laminar debido a la elevada viscosidad del jarabe, de modo que no se requiere un modelo de turbulencia adicional en la regi\'on de producto.

Las condiciones de contorno impuestas al dominio son las siguientes. En las dos entradas de la chaqueta se prescribe el caudal m\'asico total de dise\~no dividido entre las dos ramas de acuerdo con la distribuci\'on de \'areas, una temperatura de entrada de 75~\textdegree C y un perfil de velocidad desarrollado. En las dos salidas se impone presi\'on atmosf\'erica como condici\'on de salida libre. En la interfaz agua--acero y acero--glucosa se aplica continuidad de flujo de calor y temperatura. Las superficies externas del dominio de glucosa se tratan como paredes adiab\'aticas para aislar el fen\'omeno de transferencia desde la chaqueta, mientras que la cara exterior del fondo se considera aislada t\'ermicamente, equivalente al caso con aislamiento t\'ermico de dise\~no. La malla se construye con elementos prism\'aticos y tetra\'edricos refinados en la interfaz s\'olido--fluido, especialmente en la pel\'icula de glucosa pr\'oxima a la pared, donde se concentra la resistencia t\'ermica dominante.

El acoplamiento fluido-t\'ermico se resuelve de manera simult\'anea. Una vez alcanzada la convergencia, el coeficiente global promedio se obtiene a partir del flujo de calor total transferido desde el agua hacia la glucosa, el \'area de contacto de 14~m\textsuperscript{2} y la diferencia de temperatura media entre ambos fluidos. El resultado se compara directamente con el valor obtenido del modelo anal\'itico de resistencias t\'ermicas en serie.
